Como grupo, consideramos que el principal aprendizaje del trabajo fue comprender que la corrección de un circuito depende tanto de sus operaciones como de las condiciones en que recibe y conserva los datos. La interacción con los pulsadores hizo necesario relacionar una acción manual con la captura gobernada por un reloj. Esto fundamentó la separación entre sincronización, almacenamiento y cálculo, y mostró que el procedimiento de uso también forma parte de la interfaz del hardware.

La verificación por módulos permitió distinguir responsabilidades. Comprobar una suma no asegura que el usuario cargue el operando esperado ni que el reset prevalezca sobre una carga. Al ensayar primero la ALU y el sincronizador y después el sistema completo, fue posible saber qué propiedad se comprobaba en cada etapa. Los casos de carry y overflow también ayudaron a separar el patrón binario de su interpretación numérica.

Los resultados de síntesis concuerdan con la arquitectura descripta: la ALU concentra casi toda la lógica combinacional, mientras que los 31 flip-flops se explican mediante los registros de datos, los sincronizadores y el reset de arranque. La implementación cumple las restricciones internas de \SI{100}{\mega\hertz} con márgenes positivos y sin endpoints fallidos dentro del alcance definido.

Finalmente, los resultados favorables de las herramientas deben leerse junto con sus límites. Las 866 comprobaciones respaldan los casos simulados y el reporte temporal respalda las rutas restringidas; ninguno reemplaza la observación de los controles físicos. El diseño cuenta con una base consistente para generar el bitstream y completar el ensayo sobre la placa, manteniendo explícita la diferencia entre lo verificado y lo pendiente.
